\subsection{不可分解内射模的结构}

引理 1.—— 设 A 为环，a 为 A 的理想，I 为 A-模。对每个整数 \(n \geqslant 0\)，用 \(I_{n}\) 表示 I 中被 \(a^{n}\) 零化的元素组成的子模。

\begin{enumerate}
\def\labelenumi{\alph{enumi})}
\item
  假设 A-模 I 是内射的。则 \(\mathrm{A} / \mathfrak{a}^n\)-模 \(\mathrm{I}_n\) 对每个 \(n \geqslant 0\) 都是内射的。
\item
  假设环 A 是 Noether 的，且对每个 \(n \geqslant 0\)，\(\mathrm{A} / \mathfrak{a}^n\)-模 \(\mathrm{I}_n\) 都是内射的。则诸 \(\mathrm{I}_n\) 的并是一个内射 A-模。
\item
  \(\mathbf{A} / \mathfrak{a}^n\)-模 \(\mathrm{I}_n\) 同构于 \(\operatorname{Hom}_{\mathbf{A}}(\mathbf{A} / \mathfrak{a}^n, \mathbf{I})\)，后者是内射的（A, X, 第 18 页，命题 11）。
\item
  设 J 为诸 \(I_{n}\) 之并，b 为 A 的理想，\(f : b \to J\) 为一个 A-同态。剩下要证（A, X, 第 16 页，命题 10）的是：存在 J 的元素 x，使对每个 \(f(b) = bx\) 有 \(b \in b\)。由于 b 是有限生成的，存在整数 n，使 \(f(\mathfrak{b}) \subset \mathrm{I}_{n}\)，即 \(f(\mathfrak{a}^{n}\mathfrak{b}) = 0\)。根据 III, § 3, 第 1 节，命题 1 的推论 2，存在整数 \(m \geqslant n\) 使 \(a^{m} \cap b \subset a^{n}b\)，从而 \(f(\mathfrak{a}^{m} \cap \mathfrak{b}) = 0\)。于是 f 诱导出一个 A/\(a^{m}\)-线性同态，从 \(\mathfrak{b}/(\mathfrak{a}^{m} \cap \mathfrak{b})\) 到 \(I_{m}\)；由于 A/ \(a^{m}\)-模 \(I_{m}\) 是内射的，存在 \(I_{m}\) 的元素 x，使对每个 \(f(b) = bx\) 有 \(b \in b\)，由此得 b)。
\end{enumerate}

设 \(\mathfrak{a}\) 为 A 的理想；以下我们约定 \(\mathfrak{a}^n = \mathrm{A}\)（对每个整数 \(n \leqslant 0\)）。设 E 为 A-模。对每个 \(n \in \mathbf{Z}\)，用 \(\mathrm{E}_n\) 表示 E 中被 \(\mathfrak{a}^n\) 零化的元素组成的子模；设 \(\mathrm{gr}^{\mathfrak{a}}(\mathrm{E})\) 为 \(\mathbf{Z}\) 型分次 A-模，使得对每个整数 \(\mathrm{gr}^{\mathfrak{a}}(\mathrm{E})_m = \mathrm{E}_{-m+1}/\mathrm{E}_{-m}\) 有 \(m\)。模 \(\mathrm{gr}^{\mathfrak{a}}(\mathrm{E})_m\) 在 \(m \geqslant 1\) 时为零，而 \(\mathrm{gr}^{\mathfrak{a}}(\mathrm{E})_0\) 等同于 \(\mathrm{E}_1\)。设 \(\mathrm{gr}(\mathrm{A})\) 为 A 关于 \(\mathfrak{a}\)-进滤过的伴随分次环：对每个 \(\mathrm{gr}(\mathrm{A})_n = \mathfrak{a}^n/\mathfrak{a}^{n+1}\) 我们有 \(n \in \mathbf{Z}\)。设 \(n\) 与 \(m\) 为整数。对双线性映射 \((a,x) \mapsto ax\)（从 \(\mathfrak{a}^n \times \mathrm{E}_{-m+1}\) 到 \(\mathrm{E}_{-m-n+1}\)）取商，我们导出一个 A/ \(\mathfrak{a}\)-双线性映射

\[
\alpha_ {n, m}: \mathrm{gr} (\mathrm{A}) _ {n} \times \mathrm{gr} ^ {\mathfrak {a}} (\mathrm{E}) _ {m} \longrightarrow \mathrm{gr} ^ {\mathfrak {a}} (\mathrm{E}) _ {n + m},\tag{1}
\]

它在 gr \(^{a}\) (E) 上定义了一个分次 gr(A)-模结构。对每个 \(n \in \mathbf{Z}\)，由 A/a-双线性映射 \(\alpha_{n,-n} : \text{gr}(A)_n \times \text{gr}^{a}(E)_{-n} \longrightarrow E_1\) 导出一个 A/a-线性映射 \(\beta_{E,n} : \text{gr}^{a}(E)_{-n} \longrightarrow \text{Hom}_{A/\alpha}(\text{gr}(A)_n, E_1)\)；诸映射 \(\beta_{E,n}\) 是一个分次 A/a-模同态的各分量，该同态称为典范同态。

\[
\beta_ {\mathrm{E}}: \operatorname{gr} ^ {a} (\mathrm{E}) \longrightarrow \operatorname{Homgr} _ {\mathrm{A} / \mathfrak{a}} (\operatorname{gr} (\mathrm{A}), \mathrm{E} _ {1}).\tag{2}
\]

对 \(a \in \text{gr(A)}\)、\(x \in \text{gr}^{\mathfrak{a}}(\text{E})\)，\(\beta_{\text{E}}(x)(a)\) 按定义是 \(\text{gr}^{\mathfrak{a}}(\text{E})_0 = \text{E}_1\) 中元素 \(ax\)（属于 \(\text{gr}^{\mathfrak{a}}(\text{E})\)）的分量。由此可知，\(\beta_{\text{E}}\) 是 \(\text{gr(A)}\)-线性的，只要给 \(\text{Homgr}_{\text{A}/\mathfrak{a}}(\text{gr(A)}, \text{E}_1)\) 赋以由公式 \(\text{gr(A)}\)（其中 \((bf)(a) = f(ab)\) 属于 \(a, b\)，\(\text{gr(A)}\) 属于 \(f\)）所定义的 \(\text{Homgr}_{\text{A}/\mathfrak{a}}(\text{gr(A)}, \text{E}_1)\)-模结构。

命题 2.—— 设 A 为 Noether 环，a 为 A 的理想，E 为 A-模，M 为 E 的被 a 零化的子 A-模。下列条件等价：

\begin{enumerate}
\def\labelenumi{(\roman{enumi})}
\item
  E 是 M 的一个内射包；
\item
  \(\mathrm{A}/\mathfrak{a}\)-模 \(\mathrm{E}_1\) 是 \(\mathrm{A}/\mathfrak{a}\)-模 \(\mathrm{M}\) 的内射包；模 \(\mathrm{E}\) 是诸 \(\mathrm{E}_n\) 的并，且典范映射 \(\beta_{\mathrm{E}}\) 是双射。
\end{enumerate}

假设条件 (i) 满足。A/a-模 \(E_{1}\) 是内射的（引理 1, a)），且包含 M；由于 \(A/a\) 的每个含于 \(E_{1}\) 的子模都是 E 的子 A-模，故 \(E_{1}\) 是 A/a-模 M 的一个内射包。由引理 1，诸 \(E_{n}\) 之并是 E 的一个包含 M 的内射子 A-模，故等于 E。由于 E 内射，对所有 \(n \geqslant 0\) 有正合列

\[
0 \to \operatorname{Hom} _ {\mathrm{A}} (\mathrm{A} / \mathfrak {a} ^ {n}, \mathrm{E}) \longrightarrow \operatorname{Hom} _ {\mathrm{A}} (\mathrm{A} / \mathfrak {a} ^ {n + 1}, \mathrm{E}) \longrightarrow \operatorname{Hom} _ {\mathrm{A}} (\mathfrak {a} ^ {n} / \mathfrak {a} ^ {n + 1}, \mathrm{E}) \to 0;
\]

由于 \(\mathrm{Hom}_{\mathrm{A}}(\mathrm{A}/\mathfrak{a}^{m},\mathrm{E})\) 等同于 \(\mathrm{E}_{m}\)（对每个 \(m\)），且从 \(\mathrm{Hom}_{\mathrm{A}}(\mathfrak{a}^{n}/\mathfrak{a}^{n+1},\mathrm{E}_{1})\) 到 \(\mathrm{Hom}_{\mathrm{A}}(\mathfrak{a}^{n}/\mathfrak{a}^{n+1},\mathrm{E})\) 的典范单射是双射，我们推出典范同态 \(\beta_{\mathrm{E}}\) 是双射，由此得 (ii)。

假设 (ii) 满足。设 \(e: \mathrm{M} \to \mathrm{I}\) 为 M 的一个内射包。由于 I 内射，存在 A-线性映射 \(\varphi: \mathrm{E} \to \mathrm{I}\) 延拓 \(e\)。而 \(\varphi\) 把 \(\mathrm{E}_n\) 映到 \(\mathrm{I}_n\) 中（对每个 \(n\)），从而诱导出同态 \(\mathrm{gr}^{\mathfrak{a}}(\varphi):\mathrm{gr}^{\mathfrak{a}}(\mathrm{E})\to \mathrm{gr}^{\mathfrak{a}}(\mathrm{I})\) 与 \(\varphi_{1}:\mathrm{E}_{1}\rightarrow \mathrm{I}_{1}\)，并使下图交换

\[
\begin{array}{c c c} \operatorname{gr} ^ {\mathfrak {a}} (\mathrm{E}) & \xrightarrow {\beta_ {\mathrm{E}}} & \operatorname{Homgr} _ {\mathrm{A} / \mathfrak {a}} (\operatorname{gr} (\mathrm{A}), \mathrm{E} _ {1}) \\ \operatorname{gr} ^ {\mathfrak {a}} (\varphi) \Bigg \downarrow & & \Bigg \downarrow \operatorname{Homgr} (1, \varphi_ {1}) \\ \operatorname{gr} ^ {\mathfrak {a}} (\mathrm{I}) & \xrightarrow {\beta_ {\mathrm{I}}} & \operatorname{Homgr} _ {\mathrm{A} / \mathfrak {a}} (\operatorname{gr} (\mathrm{A}), \mathrm{I} _ {1}). \end{array}
\]

由于 \(E_{1}\) 与 \(I_{1}\) 都是 A/a-模 M 的内射包，同态 \(\varphi_{1}\) 是双射；由于 \(\beta_{E}\) 与 \(\beta_{I}\) 都是双射，可知 \(\mathrm{gr}^{\mathfrak{a}}(\varphi)\) 是双射。由对 n 的归纳法，这蕴含 \(\varphi\) 诱导一个从 \(E_{n}\) 到 \(I_{n}\) 上的双射（对每个 \(n \geqslant 1\)）；因此 \(\varphi\) 是双射，由此得 (i)。

引理 2.—— 设 A 为环，a 为 A 的有限生成理想，M 为每个元素都被 a 的某个幂零化的 A-模。用 A-hat 表示 A 关于 a-进拓扑的分离完备化。

\begin{enumerate}
\def\labelenumi{\alph{enumi})}
\item
  在 M 上存在唯一的 \(\widehat{\mathrm{A}}\) -模结构，它延拓给定的 A-模结构。
\item
  M 的子 A-hat-模即其子 A-模，并且 Hom \(_{\text{A}}\) (M, P) = Hom \(_{\widehat{\text{A}}}\) (M, P) 对每个 \(\widehat{\text{A}}\) -模 P 成立。
\item
  把 \(\widehat{\mathbf{A}}\) 等同于环 \(\mathrm{A} / \mathfrak{a}^n\) 的逆极限，并给 M 赋以离散拓扑。设 \(a = (a_n)\) 为 \(\widehat{\mathbf{A}}\) 的一个元素，并设 \(x\) 为 M 的一个元素。由于 \(x\) 被 \(\mathfrak{a}\) 的某个幂零化，序列 \((a_n x)\) 平稳；用 \(ax\) 表示它的极限。映射 \((a, x) \mapsto ax\) 在 M 上定义了一个 \(\widehat{\mathbf{A}}\) -模结构，它延拓给定的 A-模结构。
\end{enumerate}

反之，设 M 上给定这样一个结构；设 \(a = (a_{n})\) 为 \(\widehat{\mathrm{A}}\) 的元素，\(x\) 为 M 的元素，\(m\) 为使 \(\mathfrak{a}^{m}x = 0\) 成立的整数。对每个整数 \(n\)，\(a - a_{n}\) 属于 \(\widehat{\mathfrak{a}^{n}}\)，后者等于 \(\mathfrak{a}^{n}\widehat{\mathrm{A}}\)（III, § 2, 第 12 节，命题 16 的推论 2）；因而等式 \(ax = a_{n}x\) 对 \(n \geqslant m\) 成立，由此得唯一性断言。

\begin{enumerate}
\def\labelenumi{\alph{enumi})}
\setcounter{enumi}{1}
\tightlist
\item
  由前所述，等式 \(Ax = \widehat{A}x\) 对每个 \(x \in M\) 都成立；于是子 \(\widehat{A}\)-模（在 \(M\) 中）就是它的子 \(A\)-模。最后，设 \(u\) 为一个 \(A\)-线性同态，它把 \(M\) 映入一个 \(\widehat{A}\)-模 \(P\)。设 \(a = (a_n)\) 为 \(\widehat{A}\) 的元素，\(x\) 为 \(M\) 的元素，\(m\) 为使 \(\mathfrak{a}^m x = 0\) 成立的整数；则有 \(\mathfrak{a}^m u(x) = 0\)。由于 \(a - a_m\) 属于 \(\mathfrak{a}^m\widehat{A}\)，有
\end{enumerate}

\[
u (a x) = u \left(a _ {m} x\right) = a _ {m} u (x) = a u (x),
\]

因此 u 是 A-hat-线性的。

命题 3.—— 设 A 为 Noether 环，p 为 A 的素理想，e : A/p → I 为 A-模 A/p 的一个内射包。对每个整数 n ≥ 0，用 I \(_{n}\) 表示 I 中被 p \(^{n}\) 零化的元素组成的子模。

\begin{enumerate}
\def\labelenumi{\alph{enumi})}
\item
  A-模 I 是诸 \(I_{n}\) 的并。单射 \(\Lambda/p \to I_{1}\) 延拓为一个从 \(\kappa(\mathfrak{p})\) 到 \(I_{1}\) 上的同构；用这个同构把 \(\kappa(\mathfrak{p})\) 与 \(I_{1}\) 等同。对每个整数 \(n \geqslant 0\)，\(\Lambda/p\)-模结构在 \(I_{n+1}/I_{n}\) 上经由纯量限制得自唯一的 \(\kappa(\mathfrak{p})\)-向量空间结构；典范同态 \(\beta_{I,-n}: I_{n+1}/I_{n} \longrightarrow \operatorname{Hom}_{\mathrm{A}/\mathfrak{p}}(\mathfrak{p}^{n}/\mathfrak{p}^{n+1}, \kappa(\mathfrak{p}))\) 是 \(\kappa(\mathfrak{p})\)-向量空间之间的同构，维数有限。
\item
  在 I 上存在唯一的 \(\mathrm{A_p}\) -模结构，它诱导出 I 的 A-模结构。典范同态 \(\widehat{\mathrm{A_p}} \longrightarrow \mathrm{End}_{\mathrm{A}}(\mathrm{I})\) 是双射。
\end{enumerate}

由 A, X, 第 20 页，例 1，A/p-模 \(\kappa(\mathfrak{p})\) 是 A/p 的一个内射包。于是由命题 2，I \(_1\) 等同于 \(\kappa(\mathfrak{p})\)，I 是诸 I \(_m\) 的并，且对每个整数 \(n \geqslant 0\)，\(\beta_{1,-n}\) 是 A/p-模的同构。对 A/p 的每个非零元素 \(a\)，以 \(a\) 为比的乘法映射在 Hom \(_{\text{A/p}}(\mathfrak{p}^n/\mathfrak{p}^{n+1}, \kappa(\mathfrak{p}))\) 中可逆，因而在 I \(_{n+1}\) /I \(_n\) 中也可逆，这就证明了 a)。

设 \(s \in A - p\)。由于乘法映射 \(s_{A/p}\) 是单射，Ker \(s_I\) 在 \(A/p\) 上的迹为零，这蕴含乘法映射 \(s_I\) 是单射。于是 \(sI\) 是 I 的一个直和项（A, X, 第 19 页，推论 4），由于 I 不可分解（\(n^o\) 1，命题 1），它等于 I，从而乘法映射 \(s_I\) 是双射。因此在 I 上存在唯一的 \(A_p\) -模结构，它诱导出 I 的 A-模结构；它唯一地延拓为一个 \(\widehat{A_p}\) -模结构（引理 2）。

对每个整数 \(n\)，由典范环同态 \(\Lambda_{\mathfrak{p}} \longrightarrow \mathrm{End}_{\mathrm{A}}(\mathrm{I})\) 导出一个 A-线性映射 \(\alpha_{n}: \mathrm{A}_{\mathfrak{p}}/\mathfrak{p}^{n}\mathrm{A}_{\mathfrak{p}} \longrightarrow \mathrm{Hom}_{\mathrm{A}}(\mathrm{I}_{n},\mathrm{I})\)。考虑下行正合的交换图

\[
\begin{array}{c c c c c c c c c} 0 & \longrightarrow & \mathfrak {p} ^ {n} A _ {\mathfrak {p}} / \mathfrak {p} ^ {n + 1} A _ {\mathfrak {p}} & \longrightarrow & A _ {\mathfrak {p}} / \mathfrak {p} ^ {n + 1} A _ {\mathfrak {p}} & \longrightarrow & A _ {\mathfrak {p}} / \mathfrak {p} ^ {n} A _ {\mathfrak {p}} & \longrightarrow & 0 \\ & & \alpha_ {n + 1} ^ {\prime} \Bigg \downarrow & & \alpha_ {n + 1} \Bigg \downarrow & & \alpha_ {n} \Bigg \downarrow \\ 0 & \longrightarrow & \operatorname{Hom} _ {A} (I _ {n + 1} / I _ {n}, I _ {1}) & \longrightarrow & \operatorname{Hom} _ {A} (I _ {n + 1}, I) & \longrightarrow & \operatorname{Hom} _ {A} (I _ {n}, I) & \longrightarrow & 0 \end{array}
\]

其中 \(\alpha_{n+1}^{\prime}\) 是由 \(\alpha_{n+1}\) 诱导的同态。考虑典范的 \(\kappa(\mathfrak{p})\)-双线性映射

\[
\alpha_ {n, - n}: \mathfrak {p} ^ {n} \mathrm{A} _ {\mathfrak {p}} / \mathfrak {p} ^ {n + 1} \mathrm{A} _ {\mathfrak {p}} \times \mathrm{I} _ {n + 1} / \mathrm{I} _ {n} \longrightarrow \mathrm{I} _ {1}
\]

（公式 (1)）。与之对应的线性映射 \(I_{n+1}/I_{n} \longrightarrow \operatorname{Hom}_{\kappa(\mathfrak{p})}(\mathfrak{p}^{n}A_{\mathfrak{p}}/\mathfrak{p}^{n+1}A_{\mathfrak{p}}, I_{1})\) 在左边的等同于 \(\beta_{I,-n}\)，而在右边的是 \(\alpha_{n+1}'\)。由于 \(\beta_{I,-n}\) 由 a) 是双射，对 \(\alpha_{n+1}'\) 同样成立；于是由上面的图，对 n 用归纳法可知 \(\alpha_{n}\) 对每个 n 都是同构。由于 I 是诸 \(I_{n}\) 的并，典范映射 \(\operatorname{End}_{\mathrm{A}}(\mathrm{I}) \to \varprojlim \operatorname{Hom}_{\mathrm{A}}(\mathrm{I}_{n}, \mathrm{I})\) 是双射；从而环同态 \(\widehat{A_{p}} \to \operatorname{End}_{\mathrm{A}}(\mathrm{I})\) ——它等同于诸映射 \(\alpha_{n}\) 的逆极限——是双射。

注.—— 由前面的证明可知，\(I_{n}\) 在 \(\widehat{A_{p}}\)（相应地，在 \(A_{p}\)）中的零化子是 \(p^{n}\widehat{A_{p}}\)（相应地，\(p^{n}A_{p}\)）。因此，A-模 \(I_{n}\) 的零化子是理想 \(p^{n}A_{p}\) 在 A 中的原像，该理想有时记作 \(\mathfrak{p}^{(n)}\)，称为素理想 p 的 n 次符号幂。

推论.—— 设 J 为一个内射 A-模，使得 \(\mathrm{Ass}_{\mathrm{A}}(\mathrm{J}) = \{\mathfrak{p}\}\)。

\begin{enumerate}
\def\labelenumi{\alph{enumi})}
\item
  \(L^{\prime}\) 典范映射 \(\mathrm{J} \to \mathrm{A}_{\mathfrak{p}} \otimes_{\mathrm{A}} \mathrm{J}\) 是双射。
\item
  用 E 表示 A/p-模 Hom \(_A\) (A/p, J)。在 E 上存在唯一的 \(\kappa(p)\)-向量空间结构，它延拓 E 的 A/p-模结构；A-模 J 同构于 I \(^{([E:\kappa(p)])}\)。
\end{enumerate}

事实上，J 同构于一个 A-模 \(I^{(\mathfrak{c})}\) ，其中 c 是一个适当的基数（第 1 节，定理 1）。当 J = I 时推论由命题立得，一般情形可立即由之推出。

